\documentclass{situationist-preprint}

\title{The Past Is Elsewhere}
\msid{TSM-03-a}
\msfamily{TSM-03}
\msdiscipline{Time, Space and Modality}
\mstype{Principal manuscript}
\mskeywords{reality of the past; presentism; growing block; eternalism; truthmaking; cross-temporal relations; B-theory}

\begin{document}
\maketitle

\begin{abstract}
Does the past exist? The presentist says that only present things exist, so that Caesar and the Battle of Hastings exist no more than tomorrow's weather. We argue that the past exists: past times and their occupants are elsewhere in time, as distant places and their occupants are elsewhere in space. The argument has three premises: truth supervenes on being, some relations hold between things at different times, and no time is metaphysically privileged. The first two exclude presentism, whose surrogates for past truthmakers are either not present features of anything or presuppose a tensed fact about which representation was correct. The third excludes the growing block, which grants the past existence but only up to a privileged present edge. What remains is that, for every moment, its past and the occupants of its past exist. The thesis privileges no moment and does not assert that the future exists: in a branching structure, relative to any moment the past is a single region and the future is not.
\end{abstract}

\section{The problem}

That the Battle of Hastings took place in 1066 is true, and it is true of something: of a battle, fought by people who struck and fell. Yet the battle is not happening, and its participants are not alive. The presentist concludes that the battle does not exist, and that what makes it true that it occurred must be found in the present \citep{Prior1970,Markosian2004}. The eternalist holds that the battle exists, though not now, as the Andes exist, though not here \citep{Quine1960,Williams1951,Sider2001}. Between them, the growing-block theorist holds that past and present things exist but future things do not, so that reality grows as time passes \citep{Broad1923,Tooley1997}.

The debate has turned on three considerations. The first is truthmaking: the presentist owes an account of what present things make it true that Hastings was fought. Presentists have proposed that the world now instantiates past-directed properties, such as having been such that a battle is fought at Hastings \citep{Bigelow1996}, or that there are abstract ersatz times, maximal propositions one of which was once true \citep{Crisp2007}. \citet{Sider2001} objected that such properties are not genuinely present features of anything and that positing them is cheating; \citet{Merricks2007} replied on the presentist's behalf that truths about the past need no truthmakers. The second consideration is cross-temporal relations: I admire Caesar, and the eruption of Vesuvius caused the burial of Pompeii, and these relations seem to require relata. The third is physics: if simultaneity is relative to a frame, there is no global present to be the boundary of what exists.

We argue that the past exists, as distant places do. The argument draws on TSM-01 (time is constituted by the B-series), TSM-15 and TSM-16 (no privileged present) and TSM-12, whose branching structure lets us state the thesis without asserting that the future exists.

\section{Framework}

\begin{definition}[Moments and their pasts]\label{d:past}
Moments are ordered by the B-relation of being earlier than, and the order branches toward the later but not toward the earlier: the moments earlier than any moment $m$ form a single linearly ordered chain, the \emph{past} of $m$. A thing or event is \emph{located in} the past of $m$ if it occurs at a moment of that chain.
\end{definition}

Definition~\ref{d:past} is the branching structure of TSM-12. It contains no A-determination: nothing singles out one moment whose past is \emph{the} past.

\begin{definition}[Existence]\label{d:exist}
To exist is to be in the domain of the unrestricted quantifier. \emph{The past exists} is the thesis that, for every moment $m$, the moments in the past of $m$ and the things and events located at them exist.
\end{definition}

By MET-24, existence is a first-order property had by everything, and the quantifier in Definition~\ref{d:exist} is unrestricted. Ordinary existential claims are often restricted (Objection~\ref{o:gone}); the thesis concerns unrestricted existence.

\begin{principle}[Truth supervenes on being]\label{p:tsb}
If two possible situations differ in what is true, they differ in what exists or in which fundamental properties and relations existing things instantiate.
\end{principle}

\begin{principle}[Cross-temporal relations]\label{p:cross}
Some relations, such as causation, ancestry and admiration, hold between things located at different times; and a relation holds only between things that exist.
\end{principle}

\begin{principle}[No privileged moment]\label{p:nopriv}
No moment is metaphysically privileged: whatever ontological status one moment has in virtue of its position in time, every moment has relative to itself.
\end{principle}

\begin{assumption}[Independence of the past]\label{a:indep}
Some truths about the past of a moment are not entailed by the intrinsic features of what is located at that moment together with the laws.
\end{assumption}

Principle~\ref{p:tsb} is the weak, supervenience form of truthmaking defended by \citet{Bigelow1988} and \citet{Lewis2001}; it does not require that each truth have a particular truthmaker, only that truth not float free of being. The \enquote{fundamental} properties are those whose instantiation is a matter of how their bearers are, as opposed to properties defined by what is true at other times. Principle~\ref{p:cross} is a datum of causal and historical explanation. Principle~\ref{p:nopriv} is the conclusion of TSM-15 and TSM-16. Assumption~\ref{a:indep} is plausible on any physics in which information is lost: a coin tossed in the ancient world and lost without trace landed heads or tails, and nothing about the world now, together with the laws, need entail which.

\section{The argument}

The three positions on the past are exhaustive. Either nothing in the past of a moment exists (presentism, held absolutely or relative to each moment), or the past exists only up to a privileged moment, the present edge (the growing block), or, for every moment, its past exists (the thesis of Definition~\ref{d:exist}). We exclude the first two.

\begin{lemma}\label{l:truth}
Presentism violates Principle~\ref{p:tsb}, given Assumption~\ref{a:indep}.
\end{lemma}

\begin{proof}
By Assumption~\ref{a:indep}, there are situations $s_1$ and $s_2$ alike in the intrinsic features of everything located at a moment $m$ and in the laws, in which a truth about the past of $m$ differs: in $s_1$ the coin landed heads, in $s_2$ tails. If presentism holds at $m$, everything that exists in either situation is located at $m$, and so $s_1$ and $s_2$ do not differ in what exists or in the intrinsic features of what exists. By Principle~\ref{p:tsb}, they must differ in some fundamental property of present things that is not intrinsic. The presentist's candidates are of two kinds. A past-directed property, such as the world's having been such that the coin lands heads \citep{Bigelow1996}, is defined by what is true at another time, and so is not fundamental in the sense of Principle~\ref{p:tsb}: to posit it is to reassert the past truth as a property, not to ground it \citep{Sider2001}. An ersatz time, an abstract maximal proposition describing the world at the moment of the toss \citep{Crisp2007}, exists in both situations alike, as do the ersatz times describing heads and tails; what differs is which of them \emph{was} true, which is itself a past-tensed truth of the kind to be grounded. Either way the difference between $s_1$ and $s_2$ is not grounded in being, contrary to Principle~\ref{p:tsb}.
\end{proof}

\begin{lemma}\label{l:cross}
Presentism violates Principle~\ref{p:cross}.
\end{lemma}

\begin{proof}
The eruption of Vesuvius in 79 caused the burial of Pompeii, and a historian today admires Pliny's account of it. By Principle~\ref{p:cross}, the relations of causing and admiring hold only between existing relata. On presentism, neither the eruption nor Pliny exists, and so the relations do not hold. Since they do hold, presentism is false.
\end{proof}

\begin{lemma}\label{l:block}
The growing block violates Principle~\ref{p:nopriv}.
\end{lemma}

\begin{proof}
On the growing block, the past of the present edge exists and nothing later than the edge does. The edge therefore has an ontological status, that of being the latest existing moment, which no moment earlier than it has relative to itself: earlier moments have existing successors. This is a status a moment has in virtue of its position in time that not every moment has relative to itself, contrary to Principle~\ref{p:nopriv}.
\end{proof}

Two familiar objections confirm the diagnosis. Past people, thinking \enquote{this is the present edge} on the same evidence as ours, are mistaken, so we have no reason to think we are not among them \citep{Bourne2002,BraddonMitchell2004}; and an edge spread across space is a global present, which TSM-16 shows does not exist.

\begin{theorem}\label{t:main}
The past exists: for every moment, the moments in its past and their occupants exist.
\end{theorem}

\begin{proof}
The three positions are exhaustive. Lemmas~\ref{l:truth} and~\ref{l:cross} exclude presentism. This holds whether presentism is held absolutely or relative to each moment, since the lemmas apply at any moment at which it is held. Lemma~\ref{l:block} excludes the growing block. What remains is the thesis of Definition~\ref{d:exist}.
\end{proof}

\begin{corollary}[Elsewhere]\label{c:elsewhere}
Past things differ from present ones only in their location relative to a standpoint. To be past is to be earlier than the time of a given thought or utterance, as to be distant is to be far from its place.
\end{corollary}

\begin{proof}
By Theorem~\ref{t:main}, past things exist unrestrictedly, as present ones do, so they do not differ in existence. By TSM-15, presentness is the relation of a time to what occurs at it, and \enquote{past} is defined from it and the B-relation of being earlier. The difference between Hastings and this morning is then a difference of B-location relative to the thought that compares them, as the difference between the Andes and this room is a difference of spatial location relative to a place.
\end{proof}

\section{The past and the future}

The analogy with space is exact for the past. It is not offered for the future, and the reason must be stated carefully, because an asymmetry between past and future can look like a growing block. Theorem~\ref{t:main} is stated from the standpoint of an arbitrary moment, and so privileges none. We make no parallel claim about the future of $m$. The reason is not that later things are unreal while earlier things are real, which would make $m$ an edge. It is that, relative to $m$, there is no single region for \enquote{the future of $m$} to denote. By Definition~\ref{d:past}, the past of $m$ is one chain, common to every history through $m$; the histories through $m$ diverge after it, and on TSM-12 future contingents at $m$ are neither true nor false. Nothing in the argument for Theorem~\ref{t:main} transfers. Lemma~\ref{l:truth} needs a determinate truth about the past that the present does not entail; future contingents are not determinate truths at $m$ and so need no grounding at $m$, while future truths that are settled at $m$ are grounded in present conditions and the laws (EPI-25). Lemma~\ref{l:cross} needs relations that hold; a relation between me and a contingent future event is not yet a relation that holds.

The asymmetry is one of determinacy, not of existence relative to a privileged moment, and it holds of every moment alike. It suffices for the perdurantism of PID-11, which needs only that a person's past and present stages exist.

\section{Objections and replies}

\begin{objection}[Paraphrase]\label{o:para}
The presentist does not need past objects to state past truths. \enquote{Caesar crossed the Rubicon} means \enquote{it was the case that Caesar crosses the Rubicon}, with a tense operator in place of quantification over past times \citep{Prior1970}.
\end{objection}

\begin{reply}
The paraphrase supplies words, not a ground. What makes \enquote{it was the case that the coin lands heads} true rather than \enquote{it was the case that the coin lands tails}? Not the present arrangement of things, by Assumption~\ref{a:indep}. If the tense operator is primitive and its truth is grounded in nothing, the presentist rejects Principle~\ref{p:tsb}, which is Objection~\ref{o:merricks}. And the paraphrase leaves Lemma~\ref{l:cross} untouched: \enquote{it was the case that the eruption causes the burial} does not say that one thing caused another, and the historian's admiration has an object.
\end{reply}

\begin{objection}[Truth without truthmakers]\label{o:merricks}
Truthmaker theory is false; many truths, such as negative existentials, are not made true by anything, and truths about the past are among them \citep{Merricks2007}.
\end{objection}

\begin{reply}
Principle~\ref{p:tsb} is weaker than truthmaker maximalism. Negative existentials satisfy it: worlds that differ over whether there are unicorns differ in what exists. Truths about the past satisfy it only if the past exists, and the presentist who rejects it holds that the world could have been exactly as it is in every respect while its history differed. Some truths would then answer to nothing at all, a cost paid for a privileged present that TSM-15 shows to be no benefit.
\end{reply}

\begin{objection}[The past is gone]\label{o:gone}
Ordinary thought holds that the past is over and done with and that the dead are no more. It is a strange view on which Caesar exists.
\end{objection}

\begin{reply}
Ordinary claims of non-existence are restricted (Definition~\ref{d:exist}). That the dead are no more is true: they are not located at the time of the claim. That Caesar exists unrestrictedly is not something ordinary thought denies, any more than it denies that distant cities exist when it says that there is nothing here but fields. The air of strangeness comes from reading an unrestricted quantifier into restricted talk.
\end{reply}

\begin{objection}[Disanalogy]\label{o:dis}
Time is not like space. We can travel to distant places and return, but not to the past; we remember the past and not the future; we can affect what lies ahead and not what lies behind. The past is not merely elsewhere.
\end{objection}

\begin{reply}
These are asymmetries of causal and epistemic access, and the B-theory explains them without differences of existence: by TSM-05, the direction of time, and with it the direction of causation and of memory, derives from the low-entropy initial state of the universe. That the past is fixed relative to any moment concerns the shape of the structure (Section~4) and holds for every moment alike. None of the asymmetries requires that past things fail to exist.
\end{reply}

\begin{objection}[The growing block keeps what is right]\label{o:block}
The growing block grants the past existence and keeps the future open. It secures everything the present manuscript wants at no greater cost.
\end{objection}

\begin{reply}
It secures it at the cost of an edge, which Lemma~\ref{l:block} shows to be a privileged moment. TSM-12 obtains openness from the branching structure, in which every moment has a settled past and an open future relative to itself, without one.
\end{reply}

\section{Conclusion}

The past exists. Presentism makes truths about the past answer to nothing and cross-temporal relations hold without relata; the growing block privileges a present edge. For every moment, its past and the occupants of its past exist, and past things differ from present ones only as distant places differ from here, by location relative to a standpoint. Relative to any moment the past is a single region and the future is not, so the thesis neither privileges a moment nor asserts that the future exists.

\begin{thebibliography}{99}\small
\bibitem[Bigelow(1988)]{Bigelow1988} Bigelow, J. (1988). \emph{The Reality of Numbers: A Physicalist's Philosophy of Mathematics}. Clarendon Press.
\bibitem[Bigelow(1996)]{Bigelow1996} Bigelow, J. (1996). Presentism and properties. \emph{Philosophical Perspectives}, 10, 35--52.
\bibitem[Bourne(2002)]{Bourne2002} Bourne, C. (2002). When am I? A tense time for some tense theorists? \emph{Australasian Journal of Philosophy}, 80(3), 359--371.
\bibitem[Braddon-Mitchell(2004)]{BraddonMitchell2004} Braddon-Mitchell, D. (2004). How do we know it is now now? \emph{Analysis}, 64(3), 199--203.
\bibitem[Broad(1923)]{Broad1923} Broad, C.~D. (1923). \emph{Scientific Thought}. Kegan Paul.
\bibitem[Crisp(2007)]{Crisp2007} Crisp, T.~M. (2007). Presentism and the grounding objection. \emph{No\^us}, 41(1), 90--109.
\bibitem[Lewis(2001)]{Lewis2001} Lewis, D. (2001). Truthmaking and difference-making. \emph{No\^us}, 35(4), 602--615.
\bibitem[Markosian(2004)]{Markosian2004} Markosian, N. (2004). A defense of presentism. In D.~W. Zimmerman (ed.), \emph{Oxford Studies in Metaphysics}, Vol.~1, 47--82. Oxford University Press.
\bibitem[Merricks(2007)]{Merricks2007} Merricks, T. (2007). \emph{Truth and Ontology}. Oxford University Press.
\bibitem[Prior(1970)]{Prior1970} Prior, A.~N. (1970). The notion of the present. \emph{Studium Generale}, 23, 245--248.
\bibitem[Quine(1960)]{Quine1960} Quine, W.~V.~O. (1960). \emph{Word and Object}. MIT Press.
\bibitem[Sider(2001)]{Sider2001} Sider, T. (2001). \emph{Four-Dimensionalism: An Ontology of Persistence and Time}. Oxford University Press.
\bibitem[Tooley(1997)]{Tooley1997} Tooley, M. (1997). \emph{Time, Tense, and Causation}. Clarendon Press.
\bibitem[Williams(1951)]{Williams1951} Williams, D.~C. (1951). The myth of passage. \emph{The Journal of Philosophy}, 48(15), 457--472.
\end{thebibliography}

\end{document}
